\clearpage
\VerdierPhysicalPageAnchor{0118}%
\begingroup
\small
\setlength{\abovedisplayskip}{4pt plus 2pt minus 2pt}%
\setlength{\belowdisplayskip}{4pt plus 2pt minus 2pt}%

\VerdierSemanticAnchor{ch2-1-2-6-exactness-complete}
は完全である。$W=Z$ とすると、恒等射
$\operatorname{id}_Z\in\operatorname{Hom}_{\mathcal D}(Z,Z)$ はある射
$s\colon Z\longrightarrow Y$ から来る。よって切断が存在する。
この切断により次の可換図式が得られる。

\VerdierSemanticAnchor{ch2-1-2-6-section-diagram}
\begin{equation*}
  \begin{tikzcd}[column sep=3.5em,row sep=3.0em]
    X \arrow[r,"{\binom{\operatorname{id}}{0}}"]
      \arrow[d,"{\operatorname{id}}"']
      & X\oplus Z \arrow[r,"{(0,\operatorname{id})}"]
        \arrow[d,"{(u,s)}"]
      & Z \arrow[d,"{\operatorname{id}}"] \\
    X \arrow[r,"u"'] & Y \arrow[r,"v"'] & Z
  \end{tikzcd}\,.
\end{equation*}

\VerdierSemanticAnchor{ch2-1-2-6-hom-five-lemma}
関手 $\operatorname{Hom}_{\mathcal D}(W,\mathord\cdot)$ を適用すると、
列 (1.2.6.2) の完全性と五項補題により、射

\VerdierSemanticAnchor{ch2-1-2-6-hom-isomorphism}
\[
  \operatorname{Hom}_{\mathcal D}\bigl(W,(u,s)\bigr)\colon
  \operatorname{Hom}_{\mathcal D}(W,X\oplus Z)
  \longrightarrow \operatorname{Hom}_{\mathcal D}(W,Y)
\]
は同型である。\space したがって射 $(u,s)$ は同型である。

\VerdierSemanticAnchor{ch2-1-2-7}
\medskip
\noindent\textbf{注意 1.2.7.}\quad
系 (1.2.6) により、任意の完全関手
$\mathcal D'\xrightarrow{F}\mathcal D$ が加法的であることを示せる。
まず、三角 $0\longrightarrow0\longrightarrow0\longrightarrow0$ は
$\mathcal D'$ の完全三角である。したがって三角

\VerdierSemanticAnchor{ch2-1-2-7-zero-triangle}
\[
  F(0)\xrightarrow{\operatorname{id}_{F(0)}}F(0)
  \xrightarrow{\operatorname{id}_{F(0)}}F(0)
  \longrightarrow F(0)[1]
\]
は $\mathcal D$ の完全三角である。(1.2.2) から $F(0)=0$ が従う。
$F$ を (1.2.6.1) の形の完全三角に適用し、さらに (1.2.6) を用いると、
$F$ が加法的であることが分かる。同様の議論により、任意の
コホモロジー的関手も加法的である。

\VerdierSemanticAnchor{ch2-1-2-8}
\medskip
\noindent\textbf{定義 1.2.8.}\quad
加法圏における射 $f\colon X\longrightarrow Y$ の\emph{分解}とは、
次の可換図式をいう。

\VerdierSemanticAnchor{ch2-1-2-8-decomposition-diagram}
\begin{equation*}
  \begin{tikzcd}[ampersand replacement=\&,column sep=5.2em,row sep=3.4em]
    X \arrow[r,"f"] \arrow[d,"\wr"']
      \& Y \arrow[d,"\wr"] \\
    Z\oplus W
      \arrow[r,"{\left(\begin{smallmatrix}0&\operatorname{id}_W\\0&0\end{smallmatrix}\right)}"]
      \& W\oplus Z'
  \end{tikzcd}\,.
\end{equation*}

\VerdierSemanticAnchor{ch2-1-2-8-decomposable}
射は分解をもつとき\emph{分解可能}であるという。

\VerdierSemanticAnchor{ch2-1-2-8-splitting}
射 $f\colon X\longrightarrow Y$ の\emph{分裂}とは、次を満たす射
$g\colon Y\longrightarrow X$ をいう。
\endgroup
